\pagestyle{plain}
\begin{center}
{\LARGE Περίληψη}\\[1cm]
\end{center}

Η ραγδαία ανάπτυξη των συσκευών Διαδικτύου των Πραγμάτων (IoT) έχει δημιουργήσει νέες ευκαιρίες για την παρακολούθηση περιβαλλοντικών παραμέτρων σε πραγματικό χρόνο. Η ποιότητα αέρα, ειδικότερα, αποτελεί κρίσιμο ζήτημα σε αστικά περιβάλλοντα, όπου ρύποι όπως το μονοξείδιο του άνθρακα (CO), τα οξείδια του αζώτου (NOx) και το βενζόλιο συνιστούν μετρήσιμο κίνδυνο για τη δημόσια υγεία. Η ανάπτυξη, ωστόσο, αξιόπιστων, επεκτάσιμων και εύκολα αναπτύξιμων συστημάτων συλλογής και επεξεργασίας δεδομένων από κατανεμημένα δίκτυα αισθητήρων παραμένει σημαντική τεχνική πρόκληση. Η παρούσα διπλωματική εργασία αντιμετωπίζει αυτό ακριβώς το πρόβλημα, προτείνοντας και υλοποιώντας ένα pipeline συλλογής, μετασχηματισμού και οπτικοποίησης δεδομένων IoT αισθητήρων ποιότητας αέρα, βασισμένο σε αρχιτεκτονική microservices.

Το σύστημα που σχεδιάστηκε καλύπτει πέντε βασικές απαιτήσεις: (α) ταυτόχρονη συλλογή δεδομένων από πολλαπλούς αισθητήρες μέσω ελαφρού πρωτοκόλλου επικοινωνίας κατάλληλου για συσκευές περιορισμένων πόρων, (β) φιλτράρισμα και μετασχηματισμό ακατέργαστων μετρήσεων σε δομημένες μετρικές συμβατές με υποδομή αποθήκευσης χρονοσειρών, (γ) αποθήκευση με ρυθμιζόμενη πολιτική διατήρησης που υποστηρίζει τόσο real-time παρακολούθηση όσο και ιστορική ανάλυση, (δ) οπτικοποίηση μέσω διαδραστικών dashboards, και (ε) πλήρη αναπαραγωγιμότητα: ολόκληρο το σύστημα αναπτύσσεται από ένα μόνο αρχείο ρύθμισης, σε οποιοδήποτε περιβάλλον.

Η μέθοδος που ακολουθήθηκε βασίζεται στην αρχιτεκτονική microservices, όπου κάθε λειτουργικό στοιχείο του pipeline υλοποιείται ως ανεξάρτητη, containerized υπηρεσία. Η επιλογή αυτής της αρχιτεκτονικής έναντι της παραδοσιακής monolithic προσέγγισης αιτιολογείται από την ανάγκη για ανεξάρτητη ανάπτυξη, δοκιμή και κλιμάκωση κάθε στοιχείου. Σε IoT περιβάλλον όπου ο αριθμός αισθητήρων ενδέχεται να μεταβάλλεται δυναμικά, η δυνατότητα κλιμάκωσης μεμονωμένων υπηρεσιών χωρίς επίπτωση στις υπόλοιπες αποτελεί κρίσιμο πλεονέκτημα.

Για την υλοποίηση επιλέχθηκε συνδυασμός τεχνολογιών αντιστοιχισμένων στις επιμέρους απαιτήσεις του συστήματος. Το πρωτόκολλο MQTT, μέσω του broker Eclipse Mosquitto, εξυπηρετεί την ελαφριά publish/subscribe επικοινωνία μεταξύ αισθητήρων και επεξεργαστή δεδομένων, αποσυνδέοντας φυσικά τους παραγωγούς από τους καταναλωτές δεδομένων. Υπηρεσία controller, υλοποιημένη σε Python με FastAPI, εγγράφεται στα MQTT topics των αισθητήρων, μετασχηματίζει τα δεδομένα σε μορφή Prometheus και τα εκθέτει μέσω HTTP endpoint. Το Prometheus συλλέγει αυτές τις μετρικές μέσω pull-based μοντέλου και τις αποθηκεύει σε βελτιστοποιημένη βάση χρονοσειρών. Το Grafana αξιοποιεί το Prometheus ως πηγή δεδομένων και παρέχει διαδραστικά dashboards με υποστήριξη templating μεταβλητών και alerting. Κάθε υπηρεσία είναι πακεταρισμένη ως Docker container, ενώ η ορχήστρωση γίνεται μέσω Docker Compose.

Για τον έλεγχο του συστήματος χωρίς φυσικό υλικό, αναπτύχθηκε υπηρεσία προσομοίωσης αισθητήρων που αναπαράγει μετρήσεις από το δημόσιο σύνολο δεδομένων UCI Air Quality. Το σύνολο αυτό περιέχει 9.357 ωριαίες μετρήσεις από πολυ-αισθητήρα διάταξη αναπτυγμένη σε ιταλική πόλη, με επαληθευμένες τιμές αναφοράς για CO, NOx, NO\textsubscript{2}, βενζόλιο, θερμοκρασία και υγρασία. Κάθε προσομοιωτής λειτουργεί ως ανεξάρτητο Docker container που συμμετέχει στο pipeline ακριβώς όπως θα έκανε ένας πραγματικός αισθητήρας, ενώ η αντικατάστασή του από φυσική συσκευή δεν απαιτεί αλλαγές στο υπόλοιπο σύστημα.

Τα αποτελέσματα επιβεβαιώνουν τα οφέλη της αρχιτεκτονικής microservices για περιβάλλοντα IoT. Η προσθήκη νέου αισθητήρα απαιτεί μόνο την προσθήκη νέας υπηρεσίας στο αρχείο ορχήστρωσης, χωρίς αλλαγές στο υπόλοιπο σύστημα. Τυχόν αποτυχία μεμονωμένης υπηρεσίας δεν διακόπτει τη λειτουργία των υπολοίπων, ενώ η επανεκκίνησή της αποκαθιστά αυτόματα τη συμμετοχή της στο pipeline. Η παρακολούθηση μέσω Prometheus και Grafana παρέχει τόσο άμεση εικόνα της κατάστασης όσο και δυνατότητα ανάλυσης ιστορικών τάσεων και εντοπισμού ανωμαλιών. Ολόκληρο το σύστημα εκκινεί πλήρως με μία εντολή (\texttt{docker compose up}) σε οποιοδήποτε περιβάλλον διαθέτει Docker, επιτυγχάνοντας πλήρη αναπαραγωγιμότητα.

Μελλοντικές επεκτάσεις περιλαμβάνουν την αντικατάσταση των προσομοιωτών με πραγματικούς αισθητήρες εντός campus, τη μετάβαση σε Kubernetes για high availability και αυτόματη κλιμάκωση, την υιοθέτηση λύσης μακροπρόθεσμης αποθήκευσης μετρικών (Thanos σε συνδυασμό με S3), τη ρύθμιση ειδοποιήσεων για υπέρβαση επικίνδυνων ορίων ρύπων, καθώς και την ανάπτυξη CI/CD pipeline για αυτοματοποιημένη δοκιμή και ανάπτυξη.
